% Generated draft: requires mathematical and source review.
\documentclass[11pt]{article}
\usepackage[margin=25mm]{geometry}
\usepackage{fontspec}
\setmainfont{Latin Modern Roman}
\usepackage{amsmath,amssymb,mathtools,unicode-math}
\setmathfont{Latin Modern Math}
\usepackage{enumitem,needspace,xurl,hyperref}
\setlength{\parindent}{0pt}
\setlength{\parskip}{5pt}
\setlength{\emergencystretch}{3em}
\newcommand{\N}{\mathbb{N}}
\newcommand{\Z}{\mathbb{Z}}
\newcommand{\Q}{\mathbb{Q}}
\newcommand{\R}{\mathbb{R}}
\newcommand{\C}{\mathbb{C}}
\newcommand{\F}{\mathbb{F}}
\newcommand{\A}{\mathbb{A}}
\newcommand{\PP}{\mathbb{P}}
\newcommand{\E}{\mathbb{E}}
\newcommand{\eps}{\varepsilon}
\newcommand{\dd}{\,\mathrm{d}}
\begin{document}
\section*{UnsolvedMath ID: 800223}
\textbf{Source problem code:} catalogue entry
\subsection{Definitions and mathematical statement}
A domain is a nonempty connected open set. A domain \(D\subset\C^n\) is pseudoconvex if it admits a continuous plurisubharmonic exhaustion function \(\psi:D\to\R\), meaning that \(\{z\in D:\psi(z)<c\}\) is relatively compact in \(D\) for every \(c\in\R\). Define
\[
A(D)=\{f\in C(\overline D;\C):f|_D\text{ is holomorphic}\}.
\]
A point \(p\in\partial D\) is a peak point for \(A(D)\) if there exists \(f\in A(D)\) such that
\[
f(p)=1,\qquad |f(z)|<1\quad\text{for every }z\in\overline D\setminus\{p\}.
\]
The supplied question asks whether every boundary point of every bounded pseudoconvex domain \(D\subset\C^n\) of finite type is a peak point for \(A(D)\). The record does not specify boundary regularity or the meaning of finite type, so these hypotheses cannot be fixed unconditionally from the supplied text.

For an explicit conditional formulation, suppose that the boundary is \(C^\infty\) and that finite type means finite D'Angelo type at every boundary point. For \(p\in\partial D\), let \(r\) be a smooth local real defining function, with \(D=\{r<0\}\) locally and \(\mathrm{d}r(p)\ne0\). Define
\[
\Delta_1(\partial D,p)=\sup_{\gamma}\frac{\nu_0(r\circ\gamma)}{\nu_0(\gamma-p)},
\]
where \(\gamma:(\C,0)\to(\C^n,p)\) ranges over nonconstant holomorphic curve germs, \(\nu_0(r\circ\gamma)\) is its order of vanishing as a smooth function of the real and imaginary parts of the curve parameter, and \(\nu_0(\gamma-p)\) is the minimum vanishing order of its nonzero coordinate germs. Infinite vanishing order is allowed. Under this additional convention, the question is:
\[
\begin{gathered}
\text{For every integer }n\ge1\text{ and every bounded pseudoconvex domain }D\subset\C^n,\\
\bigl[\partial D\text{ is }C^\infty\text{ and }\Delta_1(\partial D,q)<\infty\text{ for every }q\in\partial D\bigr]\\
\Longrightarrow\quad
\forall p\in\partial D\ \exists f\in A(D):\quad
f(p)=1\ \text{ and }\ \forall z\in\overline D\setminus\{p\},\ |f(z)|<1\;?
\end{gathered}
\]
The smoothness and D'Angelo convention in this displayed formulation are explicit interpretive additions requiring source confirmation. The wording also leaves unresolved whether finite type is required on the whole boundary or only at the point under consideration; the display uses the whole-boundary reading.
\subsection{Short English statement}
The question asks whether every boundary point of a bounded pseudoconvex domain of finite type admits a holomorphic function continuous on the closure that equals \(1\) there and has modulus strictly less than \(1\) everywhere else on the closure. Boundary regularity, the intended notion of finite type, and whether its hypothesis is global or pointwise require confirmation. The supplied registry labels the problem open; that label is not independently verified.
\subsection{Sources}
\begin{itemize}
\item Local UnsolvedMath registry: permanent ID \(800223\), source problem code \texttt{catalogue entry}, title ``Several Complex Variables Problem List --- Problem 19.'' The record attributes the problem to Buzzard, Fornaess, and Krantz, with proposed year 1995, and supplies the status \texttt{open}.
\item Buzzard, Fornaess, and Krantz, \emph{The Several Complex Variables Problem List} (1995), Problem 19. Supplied source URL: \url{https://arxiv.org/abs/math/9509202}. The supplied provenance reports access on 2026-10-16 and extraction by \texttt{arxiv\_latex}; these are source-record provenance statements, not an independent visit or verification.
\end{itemize}
\end{document}
